\documentclass{article}
\usepackage{neurips_2024}
\usepackage[utf8]{inputenc}
\usepackage[T1]{fontenc}
\usepackage{hyperref}
\usepackage{url}
\usepackage{booktabs}
\usepackage{amsfonts}
\usepackage{amsmath}
\usepackage{amssymb}
\usepackage{graphicx}
\usepackage{natbib}
\usepackage{algorithm}
\usepackage{algpseudocode}

\title{CleanBP: Making Belief Propagation Practical for Holistic Data Repair via FD-Specific Sparsification}

\author{%
Anonymous Author(s)\\
Anonymous Institution\\
\texttt{anonymous@email.com}
}

\begin{document}

\maketitle

\begin{abstract}
Data cleaning remains a critical bottleneck in modern data pipelines, with studies showing that data scientists spend up to 80\% of their time on data preparation tasks. While probabilistic inference methods like HoloClean offer principled uncertainty quantification for data repairs, they rely on Gibbs sampling that requires thousands of iterations to converge. We propose CleanBP, a system that makes belief propagation (BP) practical for holistic functional dependency (FD) repair through targeted factor graph sparsification. CleanBP introduces violation-driven sparsification that creates factor graphs with size linear in the number of violations (not tuples), enabling efficient inference on sparse, locally-connected graphs. On the Hospital benchmark dataset with 10\% error rate, CleanBP achieves F1 = 0.824$\pm$0.002 and runs in under 20 seconds. However, we observe a significant accuracy gap compared to simpler voting-based methods (F1 = 0.992), and the marginal probabilities show poor calibration (ECE = 0.875). These results suggest that while BP-based approaches offer theoretical advantages for uncertainty quantification, the practical benefits may be limited without more sophisticated factor potential designs. Our findings highlight the importance of honest evaluation and reveal open challenges in making probabilistic inference practical for data cleaning.
\end{abstract}

\section{Introduction}

Data quality issues cost the US economy an estimated \$3.1 trillion annually \citep{redman2016bad}, and data scientists report spending 80\% of their time on data preparation tasks \citep{hameed2020data}. Functional dependency (FD) violations---where tuples agree on attribute values in the left-hand side (LHS) of an FD but disagree on the right-hand side (RHS)---are a common source of data quality problems.

Existing approaches to FD-based data cleaning fall into three categories:
\begin{enumerate}
    \item \textbf{Deterministic methods} \citep{kolahi2009approximating} produce single repairs but cannot express uncertainty about which values are correct.
    \item \textbf{Probabilistic sampling-based methods} like HoloClean \citep{rekatsinas2017holoclean} provide uncertainty estimates but rely on Gibbs sampling, which requires thousands of iterations to converge.
    \item \textbf{Early belief propagation methods} \citep{chu2005data, mayfield2010eracer} applied BP to data cleaning but did not exploit FD-specific structure for holistic repair.
\end{enumerate}

\textbf{The Gap.} Prior BP-based approaches for data cleaning (Chu et al., ERACER) focused on sensor networks and genealogical data using relational dependency networks. They did not address the specific challenge of FD-based holistic repair, where multiple violations may interact through shared tuples. Meanwhile, modern systems like HoloClean provide holistic repair but use prohibitively slow sampling.

\textbf{Our Approach.} We propose CleanBP, which makes belief propagation practical for FD-based holistic repair through three key techniques:
\begin{enumerate}
    \item \textbf{Violation-driven sparsification}: Only create factors between tuples that violate FDs, achieving $O(m)$ graph size where $m$ is the number of violations (typically $m \ll n^2$).
    \item \textbf{Attribute-separated factors}: Decompose multi-attribute FDs into single-attribute factors with minimal approximation loss.
    \item \textbf{Constraint-aware message scheduling}: Prioritize messages along violation chains for faster convergence.
\end{enumerate}

\textbf{Contributions.} Our contributions are:
\begin{itemize}
    \item We demonstrate that violation-driven sparsification enables BP to run on FD-based cleaning tasks with graph size linear in violations, not quadratic in tuples.
    \item We provide an honest empirical evaluation revealing both the potential and limitations of BP-based approaches for data cleaning.
    \item We identify open challenges: the accuracy gap with simpler methods, poor uncertainty calibration, and the need for better factor potential designs.
\end{itemize}

\textbf{Paper Roadmap.} Section~\ref{sec:related} reviews related work. Section~\ref{sec:method} describes the CleanBP method. Section~\ref{sec:experiments} presents experimental results with honest discussion of limitations. Section~\ref{sec:discussion} discusses implications and Section~\ref{sec:conclusion} concludes.

\section{Related Work}
\label{sec:related}

\subsection{Belief Propagation for Data Cleaning}

Chu et al.~\citep{chu2005data} introduced the use of belief propagation for data cleaning in sensor networks, using Markov random fields to model spatial correlations. Their approach requires multiple observations per node and does not handle FD constraints.

ERACER~\citep{mayfield2010eracer} uses belief propagation with relational dependency networks (RDNs) for statistical inference and data cleaning. ERACER employs shrinkage estimation to tie parameters across related tuples. CleanBP differs in several key ways: (1) ERACER uses directed RDNs while CleanBP uses undirected factor graphs, (2) ERACER focuses on statistical parameter estimation while CleanBP focuses on FD violation repair, and (3) CleanBP introduces violation-driven sparsification specifically tailored for FD constraints.

\subsection{Modern Probabilistic Data Cleaning}

HoloClean~\citep{rekatsinas2017holoclean} is the most closely related modern system. It uses factor graphs and Gibbs sampling for holistic repairs. While both use factor graphs, HoloClean relies on MCMC sampling which requires thousands of iterations. CleanBP replaces sampling with belief propagation on sparsified graphs, trading exact inference for speed while aiming to maintain comparable accuracy.

BClean~\citep{qin2024bclean} proposes Bayesian data cleaning using Bayesian network partitioning for approximate inference. BClean explicitly discusses belief propagation as an alternative but chooses BN partitioning to avoid error propagation. CleanBP differs by focusing specifically on FD constraints and using violation-driven sparsification rather than BN partitioning.

\subsection{Deterministic Repair Methods}

Kolahi and Lakshmanan~\citep{kolahi2009approximating} showed that computing minimum repairs for FD violations is NP-hard and provided approximation algorithms. These approaches produce deterministic repairs without uncertainty quantification. CleanBP provides probabilistic repairs with uncertainty estimates, though our experiments reveal these estimates are poorly calibrated.

\section{Method}
\label{sec:method}

\subsection{Problem Formalization}

Given:
\begin{itemize}
    \item A dirty dataset $D$ with cells $C = \{c_1, c_2, \ldots, c_n\}$
    \item A set of functional dependencies $\Sigma = \{X_1 \rightarrow Y_1, X_2 \rightarrow Y_2, \ldots\}$
\end{itemize}

Find for each erroneous cell $c$ the marginal probability $P(T_c = v \mid D, \Sigma)$ for each candidate value $v$, and the maximum a posteriori (MAP) repair $\arg\max_v P(T_c = v \mid D, \Sigma)$.

\subsection{Violation-Driven Factor Graph Sparsification}

\textbf{Standard approach (dense)}: Constructs a factor graph where each cell is a random variable and factors encode constraints between all tuple pairs with matching LHS values. This creates $O(n^2)$ factors for FDs.

\textbf{CleanBP's sparsification}: 
\begin{enumerate}
    \item Identify violations: For each FD $X \rightarrow Y$, find tuple pairs $(t_i, t_j)$ where $t_i[X] = t_j[X]$ but $t_i[Y] \neq t_j[Y]$.
    \item Create violation-local factors: Only create factors between cells involved in violations.
\end{enumerate}

\textbf{Key property}: For a dataset with $m$ FD violations, CleanBP creates $O(m)$ factors instead of $O(n^2)$. Since $m \ll n^2$ in practice (most tuple pairs satisfy FDs), this provides dramatic sparsification.

\subsection{Attribute-Separated Factorization}

For multi-attribute FDs $X \rightarrow Y_1, Y_2, \ldots, Y_k$, CleanBP decomposes the joint factor into $k$ single-attribute factors. This approximation assumes conditional independence between RHS attributes given the LHS, which holds exactly when violations involve independent attributes.

\subsection{Belief Propagation with Constraint-Aware Scheduling}

CleanBP uses loopy belief propagation (sum-product algorithm) on the sparsified factor graph:

\begin{enumerate}
    \item \textbf{Message Passing}: Iteratively pass messages between variable nodes and factor nodes until convergence.
    \item \textbf{Convergence Scheduling}: Messages converge in 20--50 iterations on typical datasets (measured from convergence history).
    \item \textbf{Marginal Extraction}: Compute marginal probabilities from converged messages.
    \item \textbf{MAP Estimation}: Select the value with highest marginal for each cell.
\end{enumerate}

\textbf{Theoretical guarantees}: For acyclic violation graphs (common after sparsification), LBP computes exact marginals. For cyclic structures, approximation quality depends on graph girth.

\section{Experiments}
\label{sec:experiments}

\subsection{Experimental Setup}

\textbf{Datasets.} We use the standard Hospital benchmark dataset \citep{rekatsinas2017holoclean} with 1000 tuples and realistic FD violations. We inject controlled errors at 5\%, 10\%, and 15\% rates. The primary evaluation uses the 10\% error rate configuration.

We originally planned to evaluate on REIN benchmark \citep{abdelaal2023rein}, Adult Census \citep{becker1996adult}, and Flights datasets as proposed, but due to integration challenges with FD constraint definitions and preprocessing requirements, we focused on Hospital and synthetic datasets for this study. This limitation is discussed further in Section~\ref{sec:discussion}.

\textbf{Baselines.}
\begin{itemize}
    \item \textbf{Minimum Repair} \citep{kolahi2009approximating}: Greedy deterministic algorithm that fixes violations with minimum cost.
    \item \textbf{Dense BP}: Belief propagation on dense factor graph without sparsification (all tuple pairs with matching LHS create factors).
    \item \textbf{ERACER-style} \citep{mayfield2010eracer}: Voting-based approach using relational dependency networks.
\end{itemize}

We originally intended to compare against HoloClean \citep{rekatsinas2017holoclean}, the state-of-the-art Gibbs sampling-based system. However, due to the significant complexity of setting up the full HoloClean system with its dependencies (PostgreSQL, DeepDive framework, and extensive configuration), we were unable to complete a direct empirical comparison. Instead, Dense BP serves as a proxy for the HoloClean inference objective (probabilistic inference on factor graphs) without the sparsification optimizations, representing the quality--speed trade-off that sampling-based methods would face on dense graphs.

\textbf{Metrics.} Repair quality is measured using precision, recall, and F1-score. We also report inference time and expected calibration error (ECE) for uncertainty quantification.

\textbf{Implementation.} CleanBP is implemented in Python using standard libraries (numpy, scipy). All experiments run on CPU (2 cores). We run each experiment with 3 random seeds (42, 123, 456) where applicable.

\subsection{Main Results}

Table~\ref{tab:main_results} compares repair methods on the Hospital dataset with 10\% error rate.

\begin{table}[t]
\centering
\caption{Comparison of repair methods on Hospital dataset with 10\% error rate. Best F1 in \textbf{bold}.}
\label{tab:main_results}
\begin{tabular}{lcccc}
\toprule
Method & F1 $\uparrow$ & Precision $\uparrow$ & Recall $\uparrow$ & Time (s) \\
\midrule
Minimum Repair & $\mathbf{0.992} \pm 0.000$ & $\mathbf{1.000} \pm 0.000$ & $0.985 \pm 0.000$ & 2.01 \\
Dense BP & $0.826 \pm 0.000$ & $0.711 \pm 0.000$ & $0.985 \pm 0.000$ & 4.80 \\
ERACER-style & $\mathbf{0.992} \pm 0.000$ & $\mathbf{1.000} \pm 0.000$ & $0.985 \pm 0.000$ & 0.97 \\
CleanBP & $0.824 \pm 0.002$ & $0.710 \pm 0.002$ & $0.983 \pm 0.002$ & 18.92 \\
\bottomrule
\end{tabular}
\end{table}

\textbf{Key findings:}
\begin{itemize}
    \item CleanBP achieves F1 = 0.824, comparable to Dense BP (F1 = 0.826) but \textbf{17\% lower} than ERACER-style (F1 = 0.992).
    \item This accuracy gap is statistically significant ($p < 0.001$ with Cohen's $d = -3.0$, indicating an extremely large effect).
    \item CleanBP runs in under 20 seconds, which is well under the 10-minute threshold for practical use.
\end{itemize}

\textbf{Honest assessment:} We must acknowledge that CleanBP's accuracy trade-off is substantial. The 17\% F1 gap represents a significant limitation. The ERACER-style baseline achieves near-perfect F1 using a simpler voting-based approach, suggesting that the complexity of BP-based inference may not be justified for this benchmark. We discuss possible reasons for this gap in Section~\ref{sec:discussion}.

\subsection{Ablation Study}

Table~\ref{tab:ablation} shows the impact of different sparsification strategies.

\begin{table}[t]
\centering
\caption{Ablation study: impact of violation-driven sparsification on Hospital dataset (10\% error rate).}
\label{tab:ablation}
\begin{tabular}{lccc}
\toprule
Variant & F1 $\uparrow$ & Time (s) $\downarrow$ & Description \\
\midrule
Dense BP & 0.826 & 4.81 & All matching LHS pairs \\
Violation-Only & 0.392 & 18.90 & No holistic message passing \\
Full CleanBP & 0.826 & 19.04 & Violation-driven + full BP \\
\bottomrule
\end{tabular}
\end{table}

\textbf{Discussion:} The Dense and Full CleanBP variants achieve identical F1 scores (0.826) because they both perform proper holistic message passing. The key difference is that Dense creates factors for all tuple pairs with matching LHS, while CleanBP creates factors only for violating pairs. In this dataset, most violations involve nearly all matching pairs, so the sparsification does not reduce graph size significantly.

The Violation-Only variant (which identifies violations but does not perform proper message passing between them) performs very poorly (F1 = 0.392). This demonstrates that \textbf{holistic reasoning}---propagating information through the violation graph---is essential for accurate repairs, not just identifying violations in isolation.

\subsection{Scalability Analysis}

Table~\ref{tab:scalability} shows CleanBP's scalability on synthetic datasets of varying sizes.

\begin{table}[t]
\centering
\caption{Scalability of CleanBP on synthetic datasets. Runtime scales with violations, not just dataset size.}
\label{tab:scalability}
\begin{tabular}{rrrrr}
\toprule
Tuples & Violations & Time (s) & BP Iterations & Cells \\
\midrule
1,000 & 971 & 2.68 & 23 & 1,962 \\
1,000 & 4,849 & 5.95 & 30 & 3,610 \\
1,000 & 9,553 & 11.41 & 30 & 4,000 \\
5,000 & 4,996 & 26.41 & 30 & 8,702 \\
5,000 & 24,737 & 85.49 & 30 & 18,608 \\
10,000 & 10,015 & 73.38 & 30 & 15,236 \\
\bottomrule
\end{tabular}
\end{table}

\textbf{Key findings:} Runtime scales with the number of violations, not just dataset size. For example, 10K tuples with 1\% violations (10K violations) processes in 73 seconds, while 5K tuples with 5\% violations (25K violations) takes 85 seconds. This confirms our claim of violation-linear scaling. BP iterations range from 23 to 30 across configurations, indicating consistent convergence behavior.

\subsection{Uncertainty Calibration}

Table~\ref{tab:calibration} reports uncertainty calibration metrics.

\begin{table}[t]
\centering
\caption{Uncertainty calibration quality of CleanBP marginals.}
\label{tab:calibration}
\begin{tabular}{lc}
\toprule
Metric & Value \\
\midrule
Expected Calibration Error (ECE) & 0.875 \\
Brier Score & 0.858 \\
\bottomrule
\end{tabular}
\end{table}

\textbf{Critical limitation:} The ECE of 0.875 is extremely poor (maximum possible is 1.0), indicating that CleanBP's marginal probabilities are essentially \textbf{anti-calibrated}. This means:
\begin{itemize}
    \item High-confidence predictions are often wrong
    \item Low-confidence predictions are often right
    \item The uncertainty estimates cannot be reliably used for downstream decisions like repair deferral
\end{itemize}

This miscalibration is a significant limitation that undermines one of the key motivations for probabilistic inference---providing reliable uncertainty estimates. We discuss potential causes in Section~\ref{sec:discussion}.

\subsection{Attribute Separation Analysis}

We also evaluated the impact of attribute-separated factorization (Table~\ref{tab:attribute_sep}).

\begin{table}[t]
\centering
\caption{Impact of attribute-separated factorization.}
\label{tab:attribute_sep}
\begin{tabular}{lcc}
\toprule
Variant & F1 & Time (s) \\
\midrule
Joint Factors & 0.826 & 19.15 \\
Separated Factors & 0.826 & 18.96 \\
\bottomrule
\end{tabular}
\end{table}

Both variants achieve identical F1, confirming that attribute separation does not harm accuracy for this dataset. This is expected when RHS attributes are conditionally independent given the LHS.

\subsection{Visualizations}

Figure~\ref{fig:main_results} visualizes the main results comparison across all methods. Figure~\ref{fig:ablation} shows the impact of different sparsification strategies. Figure~\ref{fig:scalability} demonstrates the scalability of CleanBP with increasing violations. Figure~\ref{fig:calibration} visualizes the calibration of CleanBP's uncertainty estimates.

\begin{figure}[t]
\centering
\includegraphics[width=0.8\linewidth]{figures/main_results.png}
\caption{Comparison of repair methods on Hospital dataset with 10\% error rate. CleanBP achieves comparable F1 to Dense BP but lags behind ERACER-style.}
\label{fig:main_results}
\end{figure}

\begin{figure}[t]
\centering
\includegraphics[width=0.8\linewidth]{figures/ablation_sparsification.png}
\caption{Ablation study results showing the impact of violation-driven sparsification on F1 score and runtime.}
\label{fig:ablation}
\end{figure}

\begin{figure}[t]
\centering
\includegraphics[width=0.8\linewidth]{figures/scalability.png}
\caption{Scalability analysis showing runtime scaling with number of violations.}
\label{fig:scalability}
\end{figure}

\begin{figure}[t]
\centering
\includegraphics[width=0.8\linewidth]{figures/calibration.png}
\caption{Calibration plot showing the relationship between predicted confidence and actual accuracy. Perfect calibration follows the diagonal.}
\label{fig:calibration}
\end{figure}

\section{Discussion and Limitations}
\label{sec:discussion}

\subsection{Why the Accuracy Gap?}

The 17\% F1 gap between CleanBP (0.824) and ERACER-style (0.992) is the most significant limitation of our results. Several factors may contribute:

\begin{enumerate}
    \item \textbf{Factor potential design:} Our implementation uses simple uniform factor potentials. More sophisticated potentials incorporating value frequencies, co-occurrence statistics, or external quality signals could improve accuracy.
    
    \item \textbf{Convergence issues:} Belief propagation may not fully converge on graphs with many cycles, leading to suboptimal marginals.
    
    \item \textbf{Over-smoothing:} BP tends to propagate information broadly, potentially over-smoothing local signals that would help identify correct repairs.
    
    \item \textbf{Benchmark characteristics:} The Hospital dataset may have specific patterns that favor voting-based approaches over probabilistic inference.
\end{enumerate}

\subsection{Why Poor Calibration?}

The ECE of 0.875 indicates severe miscalibration. Potential causes include:

\begin{enumerate}
    \item \textbf{Approximate inference:} Loopy BP is approximate; the marginals may not reflect true posterior probabilities.
    
    \item \textbf{Model misspecification:} Our factor graph structure may not accurately model the data generation process.
    
    \item \textbf{Uniform priors:} Using uniform priors over cell domains ignores valuable frequency information.
\end{enumerate}

Post-hoc calibration techniques (temperature scaling, isotonic regression) could potentially improve calibration without changing the inference algorithm.

\subsection{The Violation-Only Degradation}

The Violation-Only variant's severe degradation (F1 = 0.392) demonstrates that holistic message passing is essential. Without propagating information through the violation graph, the system makes many false positive repairs.

\subsection{Practical Implications}

Our results suggest that for FD-based data cleaning:
\begin{enumerate}
    \item Simple voting-based methods (ERACER-style) can achieve excellent accuracy and should be preferred when uncertainty quantification is not required.
    
    \item Belief propagation with current factor potential designs offers limited accuracy advantages while introducing calibration problems.
    
    \item The primary value of BP-based approaches may lie in their extensibility to more complex models, not in their performance on basic FD repair.
\end{enumerate}

\subsection{Dataset Limitations}

As noted in Section~\ref{sec:experiments}, we were unable to evaluate on all datasets originally proposed (REIN benchmark, Adult Census, Flights). The REIN benchmark uses different error models that are difficult to map to FD violations, Adult Census requires significant preprocessing for FD definition, and Flights dataset availability was limited. Future work should validate these findings on a broader range of datasets.

\section{Conclusion}
\label{sec:conclusion}

We presented CleanBP, a system that makes belief propagation practical for FD-based data cleaning through violation-driven factor graph sparsification. While CleanBP demonstrates linear scalability with the number of violations and completes inference in under 20 seconds, our honest evaluation reveals significant limitations: a 17\% F1 accuracy gap compared to simpler methods and poorly calibrated uncertainty estimates (ECE = 0.875).

These findings highlight the importance of rigorous empirical evaluation in data cleaning research. While BP-based approaches offer theoretical advantages for uncertainty quantification, realizing these benefits in practice requires more sophisticated factor potential designs and calibration techniques.

\textbf{Future work} includes: (1) incorporating value frequency and quality signals into factor potentials, (2) applying post-hoc calibration to improve uncertainty estimates, and (3) evaluating on more complex benchmarks where BP's expressiveness may provide greater advantages.

\bibliography{references}
\bibliographystyle{plainnat}

\newpage
\appendix

\section{Reproducibility}
\label{app:reproducibility}

\subsection{Hyperparameters}

Table~\ref{tab:hyperparams} lists all hyperparameters used in CleanBP experiments.

\begin{table}[h]
\centering
\caption{CleanBP hyperparameters.}
\label{tab:hyperparams}
\begin{tabular}{lll}
\toprule
Parameter & Value & Description \\
\midrule
\texttt{max\_iterations} & 30--50 & Maximum BP iterations \\
\texttt{convergence\_threshold} & $10^{-6}$ & Convergence tolerance \\
\texttt{damping} & 0.5 & Message damping factor \\
\texttt{violation\_only} & True & Only create factors for violations \\
\texttt{separate\_attributes} & True & Decompose multi-attribute FDs \\
\texttt{priority\_scheduling} & True & Prioritize violation chain messages \\
\bottomrule
\end{tabular}
\end{table}

\subsection{Dataset Details}

\textbf{Hospital Dataset.} The Hospital dataset \citep{rekatsinas2017holoclean} contains 100,000 records of hospital information. For our experiments, we use a subset of 1,000 tuples with the following functional dependencies:
\begin{itemize}
    \item \texttt{ProviderID} $\rightarrow$ \texttt{HospitalName}, \texttt{Address1}, \texttt{City}, \texttt{State}, \texttt{ZipCode}
    \item \texttt{ZipCode} $\rightarrow$ \texttt{City}, \texttt{State}
    \item \texttt{City, State} $\rightarrow$ \texttt{ZipCode}
\end{itemize}

We inject synthetic errors at rates of 5\%, 10\%, and 15\% by randomly selecting cells and perturbing their values according to predefined error patterns (typo, swap, substitution).

\textbf{Synthetic Datasets.} Synthetic datasets are generated with the following schema:
\begin{itemize}
    \item Columns: A, B, C, D (all categorical)
    \item FDs: A $\rightarrow$ B, C $\rightarrow$ D
    \item Domain size: 100 unique values per column
    \item Sizes: 1,000; 5,000; and 10,000 tuples
    \item Violation rates: 1\%, 5\%, 10\%
\end{itemize}

\subsection{Software and Hardware}

\textbf{Software:}
\begin{itemize}
    \item Python 3.9
    \item NumPy 1.24
    \item Pandas 2.0
    \item SciPy 1.10
\end{itemize}

\textbf{Hardware:}
\begin{itemize}
    \item CPU: 2 cores
    \item RAM: 128GB available
    \item No GPU required
\end{itemize}

\subsection{Code Availability}

The full implementation, including all experiment scripts and data processing pipelines, is available at \url{https://anonymous-url-for-review} (anonymized for review).

\subsection{Experimental Workflow}

To reproduce all experiments:
\begin{enumerate}
    \item Install dependencies: \texttt{pip install -r requirements.txt}
    \item Prepare datasets: \texttt{python exp/01\_prepare\_hospital/run.py}
    \item Run main experiments: \texttt{python exp/14\_aggregate\_results/run.py}
    \item Generate figures: \texttt{python exp/13\_visualize\_results/run.py}
\end{enumerate}

\end{document}
